\section{Three and four rows}\label{sec:narrow}

In this section coordinates start at zero: columns are $0,\ldots,m-1$ and rows
are $0,\ldots,n-1$.

The plan is as follows. Most rectangles are axis-aligned, and for those the
question turns into a question about a graph: Lemma~\ref{lem:girth} says
exactly which blackouts keep the axis-aligned rectangles apart. A classical
extremal bound on that graph gives the upper bound (Theorem~\ref{thm:rowbound}).
On three and four rows there are only a few kinds of tilted rectangles
(Lemma~\ref{lem:tilted}), so for the lower bound we can write down a blackout
and check the tilted rectangles by hand.

To see where the graph comes from, think of a black point $(x,y)$ as a link
between row $y$ and column $x$. An axis-aligned rectangle uses 2 rows and 2
columns, and its 4 corners are the 4 possible links between them. If all 4 are
black, the 4 links form a closed loop of length $4$, and the next lemma shows
that such loops, and loops of length $6$, are exactly what causes confusion.

\begin{lemma}[axis-aligned identification and short cycles]\label{lem:girth}
Suppose $m\ge3$. Form a bipartite graph $G_S$ whose vertices are the rows and
columns, with an edge between row $y$ and column $x$ exactly when $(x,y)\in S$.
The axis-aligned rectangles have distinct presentations if and only if $G_S$
has neither a cycle of length four nor a cycle of length six.
\end{lemma}
\begin{proof}
If rows $a,b$ are black in both columns $x,z$, choose a third column $w$.
The rectangles on these rows and column pairs $\{x,w\}$ and $\{z,w\}$
have the same presentation. Thus a four-cycle is forbidden.
If a six-cycle has successive vertices
$a,x,b,y,c,z,a$, the rectangles with row/column pairs
$(\{a,b\},\{x,y\})$ and $(\{a,c\},\{y,z\})$ have every corner black
except possibly their common corner $(y,a)$. Their presentations coincide.

Conversely, assume both cycles are absent. Let $R_1\ne R_2$ be axis-aligned
rectangles with the same presentation $P$; we derive a contradiction by
cases on $|P|$. If $|P|\ge3$, Lemma~\ref{lem:threecorners} gives $R_1=R_2$.
If $|P|=0$, then $R_1$ is entirely black, and its 4 corners are the 4
edges of a four-cycle. Suppose $|P|=2$. If the 2 visible points are in different rows and
different columns, they are opposite corners of both rectangles, and 2
opposite corners determine an axis-aligned rectangle. If they lie in one row
$a$, at columns $x$ and $y$, then $R_1$ and $R_2$ use rows $\{a,b\}$ and
$\{a,c\}$ with $b\ne c$, and the 4 points $(x,b),(y,b),(x,c),(y,c)$ are black, a
four-cycle. The case of one column is the same with rows and columns exchanged. Finally, suppose $|P|=1$, so the common presentation is the single corner $(x,a)$.
Write the other row and column of the two rectangles as $(b,y)$ and $(c,z)$.
If $b=c$ or $y=z$, their black corners contain a four-cycle. Otherwise the six
black edges form the cycle $a,y,b,x,c,z,a$. Every possibility contradicts the
assumption.
\end{proof}

\begin{theorem}[general row bound]\label{thm:rowbound}
For $m\ge3$, every valid blackout on an $n\times m$ grid satisfies
\[
|S|\le m+\left\lfloor\frac{n^2}{4}\right\rfloor.
\]
If also $n\ge3$, the transposed bound holds as well.
\end{theorem}
\begin{proof}
For each nonempty column let $B_x$ be its set of black rows, and choose any
tree on $B_x$ (no edges if $|B_x|=1$). Superimpose these trees into a graph $H$
on the $n$ row vertices. It is simple: repeating an edge would mean two columns
share two black rows, a four-cycle in $G_S$. It is triangle-free. A triangle
cannot lie within one tree; if two edges come from one column and the third
from another, those columns share two rows; if the three edges come from three
columns, they give a six-cycle in $G_S$. All cases are forbidden by
Lemma~\ref{lem:girth}.

Recall that a triangle-free graph on $n$ vertices with $e>0$ edges
satisfies $d(u)+d(v)\le n$ for every edge $uv$. Summing over edges and using
Cauchy--Schwarz gives
\[
\frac{4e^2}{n}\le\sum_v d(v)^2\le ne,
\]
so $e\le\lfloor n^2/4\rfloor$, the usual Mantel bound. If $t\le m$ columns
are nonempty, the construction gives
$e=\sum_{B_x\ne\varnothing}(|B_x|-1)=|S|-t$.
Therefore $|S|=t+e\le m+\lfloor n^2/4\rfloor$.
\end{proof}

\begin{lemma}[tilted rectangles on narrow grids]\label{lem:tilted}
Every tilted rectangle on three rows is a translate of
\[
D(x,y)=\{(x,y+1),(x+1,y),(x+2,y+1),(x+1,y+2)\},
\qquad y=0.
\]
On four rows, these diamonds have $y=0,1$, and the remaining tilted rectangles
are the following four types, for $0\le x\le m-4$:
\begin{align*}
A_x&=\{(x,2),(x+1,3),(x+2,0),(x+3,1)\},\\
B_x&=\{(x,2),(x+2,3),(x+1,0),(x+3,1)\},\\
C_x&=\{(x,1),(x+1,3),(x+2,0),(x+3,2)\},\\
E_x&=\{(x,1),(x+2,3),(x+1,0),(x+3,2)\}.
\end{align*}
\end{lemma}
\begin{proof}
A tilted rectangle has a unique leftmost corner. Indeed, suppose 2 corners
have the smallest $x$-coordinate $x_0$. They are not adjacent, since no side is
vertical, so they are opposite and the midpoint of the rectangle has
$x$-coordinate $x_0$. The other 2 corners have the same midpoint, so their
$x$-coordinates add up to $2x_0$; both are at least $x_0$, so both equal $x_0$,
and all 4 corners lie on one vertical line, which is impossible.
Both sides at the leftmost corner point strictly to the right, and
since they are orthogonal one points up and one points down.
At the leftmost corner, write the side vectors as $(a,b)$ and $(c,-d)$,
where $a,b,c,d$ are positive integers. Orthogonality says $ac=bd$, and the
vertical span is $b+d$. If this span is at most two, all four integers are one.
If it is three, $\{b,d\}=\{1,2\}$ and $\{a,c\}=\{1,2\}$, giving the four
listed cases. These exhaust the possibilities.
\end{proof}

\begin{theorem}[Three-Row Theorem]\label{thm:threerow}
For every $m\ge3$, the maximum valid blackout on a $3\times m$ grid is $m+2$.
\end{theorem}
\begin{proof}
The upper bound is Theorem~\ref{thm:rowbound}. For the lower bound take
\[
S=\{(x,0):0\le x<m\}\cup\{(0,1),(0,2)\}.
\]
Its incidence graph is a tree, so its axis-aligned presentations are distinct.
By Lemma~\ref{lem:tilted}, each tilted rectangle is a diamond $D(x,0)$.
For $x>0$ it has only one black corner, hence at least three visible corners
and cannot be confused with any other rectangle. For $x=0$ its presentation
is $\{(2,1),(1,2)\}$. The unique axis-aligned rectangle with these opposite
corners also contains the white points $(1,1),(2,2)$, so cannot have this
presentation. There is no second tilted rectangle with at most two visible
corners. Thus $S$ is valid and has $m+2$ points.
\end{proof}

\begin{theorem}[Four-Row Theorem]\label{thm:fourrow}
For every $m\ge6$, the maximum valid blackout on a $4\times m$ grid is $m+4$.
\end{theorem}
\begin{proof}
The upper bound is Theorem~\ref{thm:rowbound}. For the matching construction set
\begin{align*}
S={}&\{(x,0):1\le x\le m-2\}\\
&{}\cup\{(m-2,1),(m-1,1),(0,2),(1,2),(0,3),(m-1,3)\}.
\end{align*}
Its size is $m+4$. Its columns with two black points are
$0,1,m-2,m-1$, joining row pairs
$\{2,3\},\{0,2\},\{0,1\},\{1,3\}$, respectively. Thus $G_S$ is an
eight-cycle with pendant edges, and all axis-aligned presentations are distinct.

It remains to check tilted rectangles. Any rectangle with at least three visible
corners is uniquely determined, so only rectangles with at least two black
corners need attention. All diamonds $D(x,1)$, and all $C_x,E_x$, have at most
one black corner. Put $L=\{0,1\}$ and $R=\{m-5,m-4\}$. The remaining
presentations with at most two visible corners are exactly those in the table.
In the last three rows $F=A$ or $B$, with $s=1$ or $2$, respectively.
\begin{center}\small
\begin{tabular}{lll}
\toprule
rectangle & range & visible corners\\
\midrule
$D(0,0)$ & & $\{(0,1),(2,1)\}$\\
$D(x,0)$ & $x=m-4,m-3$ & $\{(x,1),(x+1,2)\}$\\
$F_x$ & $x\in L\setminus R$ & $\{(x+s,3),(x+3,1)\}$\\
$F_x$ & $x\in R\setminus L$ & $\{(x,2),(x+s,3)\}$\\
$F_x$ & $x\in L\cap R$ & $\{(x+s,3)\}$\\
\bottomrule
\end{tabular}
\end{center}
The table follows by checking membership in $S$. Its two-point
presentations are distinct: their row pairs and horizontal separations distinguish
the types, and within a type the coordinates determine $x$. The only one-point
presentations occur at $m=6,x=1$, and are $(2,3)$ and $(3,3)$, which are distinct.

None equals an axis-aligned presentation. For $\{(0,1),(2,1)\}$ no other row
is black in both columns $0,2$. For $\{(x,1),(x+1,2)\}$ the other axis-aligned
corner $(x,2)$ is white, since $x\ge2$. In the $L\setminus R$ case,
$(x+s,1)$ is white, since $x+s\le3<m-2$. In the $R\setminus L$ case,
$(x,3)$ is white, since $0<x<m-1$. Finally, when $m=6$, an axis-aligned
rectangle with sole visible corner $(z,3)$, $z=2$ or $3$, would need another
column $c=0$ or $5$ and another row black in both columns $z,c$. No such row
exists: row zero is white at $0,5$, and rows one and two are white at $z$.
Hence $S$ is valid.
\end{proof}

\begin{lemma}[four-row upper-bound equality cases]\label{lem:fourclassification}
If $m\ge4$ and a valid blackout has $m+4$ black points on four rows, then every
column is nonempty, exactly four columns have two black points, the others have
one, and the four double columns form a four-cycle on the row vertices.
\end{lemma}
\begin{proof}
Equality in Theorem~\ref{thm:rowbound} requires every column nonempty and total
excess $\sum_x(|B_x|-1)=4$. A column containing all four rows leaves all other
columns with at most one black point, giving excess three. If a column contains
three rows, any other nonsingleton column must pair the fourth row with one of
those three. There can be at most one such column: two with the same row pair
give a four-cycle in $G_S$, and two with different row pairs give a six-cycle.
The total excess is then at most $2+1=3$. Thus all nonsingleton columns have
size two. They give a simple triangle-free graph with four vertices and four
edges (the graph $H$ in the proof of Theorem~\ref{thm:rowbound}). Such a graph is
a four-cycle: its complement has 2 edges, and if they shared a vertex $v$, the
other 3 vertices would form a triangle; so the 2 missing edges are disjoint,
and what remains is a four-cycle.
\end{proof}

\begin{corollary}[finite boundary cases, computer-assisted]\label{cor:fourboundary}
The maximum blackout is $7$ on $4\times4$ and $8$ on $4\times5$.
\end{corollary}
\begin{proof}
The upper bound leaves only a possible value $m+4$ to exclude. By
Lemma~\ref{lem:fourclassification}, candidates are generated by choosing one of
the three four-cycles on the row labels, four columns, a bijection from their
four row-pair edges to those columns, and a row for each remaining singleton
column. This is the complete list of
\[
3\cdot4!\binom m4 4^{m-4}
\]
candidates: $72$ when $m=4$ and $1440$ when $m=5$. The standard-library script
\texttt{verify\_narrow.py} checks every candidate directly against all rectangles,
generated independently from equal-length diagonals with a common midpoint.
Every candidate has a repeated presentation. The same script checks that
the blackout consisting of row zero and column zero is valid in these two grids,
attaining $m+3$.
\end{proof}

The verifier also checks Lemma~\ref{lem:girth} on all 74,304 blackouts of five
small grids, compares the lists of tilted rectangles in Lemma~\ref{lem:tilted}
with the diagonal enumeration, and checks the constructions for widths
$3,\ldots,40,64,100,128$ where they apply. These checks are only a test of the
proofs above; the proofs do not depend on them.
